\chapter{标准库能力地图}

\begin{longtable}{p{0.24\linewidth}p{0.30\linewidth}p{0.36\linewidth}}
\toprule
领域 & 常用设施 & 学习重点 \\
\midrule
字符串 & \code{std::string}, \code{std::string_view} & 拥有和观察的区别 \\
连续容器 & \code{std::vector}, \code{std::array}, \code{std::span} & 生命周期、扩容、边界 \\
关联容器 & \code{std::map}, \code{std::unordered_map} & 有序性、哈希、复杂度 \\
资源管理 & \code{unique_ptr}, \code{shared_ptr}, RAII 类型 & 所有权表达 \\
错误表达 & \code{optional}, 异常, \code{error_code} & 错误是否进入接口 \\
算法 & \code{find}, \code{sort}, \code{copy_if}, \code{accumulate} & 用标准词汇表达意图 \\
泛型 & template, concept, type traits & 能力约束和实例化成本 \\
并发 & \code{thread}, \code{mutex}, \code{atomic}, \code{condition_variable} & 共享状态和同步边界 \\
\bottomrule
\end{longtable}

这张表不是背诵列表。每学一个设施，都要问三件事：它是否拥有资源，它的复杂度合同是什么，它失效或出错时会怎样。

\section{按问题选择设施}

\begin{longtable}{p{0.30\linewidth}p{0.30\linewidth}p{0.30\linewidth}}
\toprule
你遇到的问题 & 优先考虑 & 先问清楚 \\
\midrule
函数可能没有结果 & \code{std::optional} & 没有结果是否还需要原因 \\
函数要返回成功或错误 & 结果结构、\code{expected} 风格类型 & 调用者是否要按错误类型分支 \\
只观察一段连续数据 & \code{std::span<const T>} & 底层数据是否活得足够久 \\
只观察字符串内容 & \code{std::string_view} & 是否会保存到调用之后 \\
独占动态对象 & \code{std::unique_ptr} & 所有权在哪里转移 \\
共享生命周期对象 & \code{std::shared_ptr}, \code{std::weak_ptr} & 是否存在循环引用 \\
按 key 快速查找 & \code{std::unordered_map} & 是否需要稳定顺序和自定义 hash \\
表达筛选转换 & \code{std::ranges}, \code{views} & 视图是否可能悬空 \\
保护共享状态 & \code{std::mutex}, RAII 锁 & 锁保护的不变量是什么 \\
\bottomrule
\end{longtable}

\section{容易误用的设施}

\begin{itemize}
  \item \code{std::string_view}：它不拥有字符，不能替代需要长期保存的 \code{std::string}。
  \item \code{std::shared_ptr}：它不是“不想思考所有权”的默认答案。
  \item \code{std::vector::reserve}：它只预留容量，不创建元素。
  \item \code{std::remove}：它不删除容器元素，需要配合 \code{erase}。
  \item ranges 视图：它常常惰性且非拥有，返回视图前要检查底层生命周期。
  \item \code{std::atomic}：它只解决特定共享变量的原子访问，不自动保护多个字段的不变量。
\end{itemize}
